\chapter{Stand der Technik / Verwandte Arbeiten}


Dieses Kapitel beschreibt den aktuellen Forschungs- und Entwicklungsstand zum Thema der Arbeit. Es stellt die wichtigsten bestehenden Ansätze, Verfahren und Systeme vor und vergleicht deren Vor- und Nachteile. Ziel ist es, die Arbeit in den Kontext bestehender Forschung einzuordnen und die eigene Herangehensweise abzugrenzen.

\begin{itemize}
\item Beschreiben Sie relevante Arbeiten chronologisch oder thematisch gruppiert (z. B. klassische Verfahren vs. aktuelle Ansätze mit Deep Learning).
\item Arbeiten Sie heraus, welche Methoden etabliert sind und welche Grenzen sie haben.
\item Zeigen Sie auf, welche Forschungslücke Ihre Arbeit schließt.
\end{itemize}

\textbf{Hinweis zur Zitierweise (IEEE-Stil):}
\begin{itemize}
\item Im IEEE-Stil werden Quellen fortlaufend nummeriert, z. B. \texttt{[1]}, \texttt{[2]}.
\item Im Text erscheinen Zitate in eckigen Klammern: „... wie in [3] gezeigt wird“.
\item Mehrere Quellen gleichzeitig: \texttt{[3]--[5]} oder \texttt{[3], [7], [9]}.
\item Die Bibliographie ist numerisch sortiert und enthält vollständige Angaben: Autor(en), Titel, Zeitschrift/Konferenz, Band, Jahr, Seiten.
\end{itemize}

\textit{Beispiel:}
\begin{quote}
Die Methode nach Smith et al. [4] erzielt bei Sprachklassifikation eine Genauigkeit von über 95%, ist jedoch rechenintensiv. Im Gegensatz dazu basiert der Ansatz von Müller [5] auf einer linearen Approximation und ist effizienter.
\end{quote}


Die Quellenangabe erfolgt dann z. B. über \texttt{\cite{knuth1984texbook}}.




\chapter{Theoretische Grundlagen}\label{ch:theorie}

Dieses Kapitel enthält die theoretischen und methodischen Grundlagen, die für das Verständnis der Arbeit erforderlich sind.
Dazu gehören Definitionen, mathematische Modelle, Fachbegriffe, Standards oder relevante Theorien.

Je nach Fachrichtung können hier behandelt werden:

\begin{itemize}
\item Mathematische oder physikalische Grundlagen, die später verwendet werden.
\item Beschreibung bestehender Verfahren, Algorithmen oder Systeme.
\item Darstellung relevanter wissenschaftlicher Arbeiten (Stand der Forschung).
\end{itemize}

Ziel ist es, die Lesenden in die Lage zu versetzen, die Methodik der Arbeit vollständig nachzuvollziehen.



